\subsection{No exhaustive filtered numerical lift of the actual cochains}

\begin{theorem}[A-026: \status{Refuted}; the full filtered numerical lift]
\label{sa08:A026}
Put \(R=\operatorname{Num}(\Z T)=\operatorname{Int}(\Z)\), with its
polynomial-degree filtration \(R_{\le w}\), and identify
\(\operatorname{gr}R=S=\Gamma_{\Z}(T)\). Write \(I_i\subset S\) for
the tail with basis \(T^{[n]}\), \(n\ge i\).
\begin{enumerate}
\item Let \(M_1,M_2\) be \(R\)-modules with exhaustive increasing
filtrations \(F_w\), with \(F_{-1}=0\) and
\(R_{\le a}F_bM_i\subset F_{a+b}M_i\). Suppose that
\(\operatorname{gr}M_i\cong I_i\) as graded \(S\)-modules, for
\(i=1,2\). Then every filtration-preserving \(R\)-bilinear map
\[
 b:M_1\otimes_R M_1\longrightarrow M_2
\]
is zero. In particular the multiplication
\(I_1\otimes_S I_1\to I_2\) cannot lift to such modules.

\item There is no contravariantly natural filtered numerical
cochain algebra \(B(X)\), for all simplicial sets \(X\), with all
the following properties. Its filtration is indexed by \(w\ge0\),
starts at zero in weight \(-1\), and is exhaustive:
\[
 B(X)\simeq\operatorname*{colim}_{w\ge0}F_wB(X).
\]
It has filtered coefficient maps
\(R_{\le w}\to F_wB(X)\), and compatible unital products
\[
 F_aB(X)\otimes^{\mathbf L}_{\Z}F_bB(X)
       \longrightarrow F_{a+b}B(X).
\]
The resulting cohomology is an \(R\)-algebra and these products
are \(R\)-bilinear. Finally there is a natural equivalence
\[
 \operatorname{gr}B(X)\simeq C_{\mathrm{PD},\mathrm{gr}}(X)
 \tag{A-26a}
 \]
as graded \(S\)-algebra objects, compatible with coefficients,
pullbacks, and the weighted cup product.
Here the filtration stages and equivalence may be taken in the
derived category, or in an enhanced category of complexes;
an actual filtered \(E_\infty\) \(R\)-algebra is included.
Only the values on the circle and the torus, with their
coefficient and projection maps, are needed for the contradiction.
\end{enumerate}
This excludes an associated-graded lift of the actual weighted
cochains. It does not exclude the weaker simplexwise total-derived
or shadow construction of Statement~\ref{sa08:D026}.
\end{theorem}

\begin{proof}[Proof of the filtered-module assertion]
Since the successive filtration quotients of \(M_i\) are
\(\Z\) in weights \(n\ge i\) and zero below \(i\), choose an
integral filtered basis \(m_n^{(i)}\), \(n\ge i\), representing
the specified elements \(T^{[n]}\). Thus
\[
 F_wM_i=\bigoplus_{n=i}^{w}\Z m_n^{(i)}.
 \tag{A-26b}
\]
Rationally, multiplication by \(m_i^{(i)}\) gives an isomorphism
\(\Q[T]\to M_i\otimes\Q\), with filtration shifted by \(i\).
Indeed its associated-graded map is an isomorphism: all the
successive coefficients of multiplication by \(T\) are
nonzero rational numbers. Induction on a finite filtration
stage proves surjectivity, and a nonzero top coefficient
proves injectivity.

Use this isomorphism to set \(m_i^{(i)}=1\) and regard all the
basis elements as rational polynomials. The graded \(S\)-action
then gives
\[
 \deg m_n^{(i)}=n-i,\qquad
 \operatorname{lc}(m_n^{(i)})=\frac{i!}{n!},
 \qquad
 Tm_n^{(i)}=(n+1)m_{n+1}^{(i)}
                 +\sum_{j=i}^{n}a_{n,j}m_j^{(i)},
 \quad a_{n,j}\in\Z.
 \tag{A-26c}
\]
The last formula proves the essential coefficient bound
\[
             m_n^{(i)}\in\frac{i!}{n!}\Z[T].
 \tag{A-26d}
\]
In fact the bound is true at \(n=i\). If true through \(n\),
all the polynomials on the left and in the lower sum in
(A-26c) have coefficients in \((i!/n!)\Z\), since
\(j!\) divides \(n!\) for \(j\le n\). Division by \(n+1\)
proves the bound at \(n+1\).

After rationalization, a bilinear \(b\) is multiplication by
some polynomial in \(\Q[T]\). Its value on the two bottom
generators lies in \(F_2M_2=\Z m_2^{(2)}\). It is consequently
multiplication by an integer \(c\). If \(c=0\), the map is zero
before rationalization as well, because \(M_2\) is torsion free.
Suppose therefore that \(c\ne0\), and choose an odd prime
\(p\) not dividing \(c\).

Localize the displayed lattices at \(p\). Set
\[
 P=T^p-T,\qquad
 f=\overline{p\,m_p^{(1)}},\qquad
 g=\overline{p\,m_p^{(2)}},\qquad
 h=\overline{p^2m_{2p}^{(2)}}\quad\text{in }\F_p[T].
 \tag{A-26e}
\]
The bars mean coefficientwise reduction of integral polynomials,
as guaranteed by (A-26d). Their degrees are respectively
\(p-1,p-2,2p-2\), and their leading coefficients are nonzero.

For \(M_1\), the reductions of
\(pF_{p+1}(M_1\otimes\Z_{(p)})\) span exactly
\(\F_pf+\F_pTf\). All basis elements below weight \(p\)
give zero on reduction after multiplication by \(p\).
The recurrence (A-26c) at \(n=p\), with its unit coefficient
\(p+1\), expresses the remaining generator
\(\overline{p\,m_{p+1}^{(1)}}\) as a nonzero multiple of
\(Tf\) plus a multiple of \(f\).
On the other hand,
\(\beta_p(T)=\binom Tp\) times the bottom generator belongs
to \(F_{p+1}M_1\), and
\[
 \overline{p\,\beta_p(T)}=\frac{P(T)}{(p-1)!}.
\]
It follows that \(P\) is a linear polynomial times \(f\).
Since \(P\) has the distinct roots \(\F_p\), there are
\(a\in\F_p\) and \(u\in\F_p^\times\) such that
\[
                 f=u\,\frac{P(T)}{T-a},
                 \qquad f(a)\ne0.
 \tag{A-26f}
\]

For \(M_2\), all basis elements of weight less than \(2p\)
have at most one factor \(p\) in their coefficient denominators.
Thus the reduction of
\(p^2F_{2p}(M_2\otimes\Z_{(p)})\) is exactly the line
\(\F_ph\). In the integral filtered basis, the product
\(\beta_p m_p^{(2)}\) has top coefficient
\(\binom{2p}{p}\), by comparison of leading polynomial
coefficients. Multiplication by \(p^2\) and reduction gives
\[
          \binom{2p}{p}h=\frac{Pg}{(p-1)!}.
 \tag{A-26g}
\]
The binomial coefficient is \(2\) modulo \(p\), hence is a unit.
For instance this follows by comparing the coefficient of
\(x^p\) in \((1+x)^{2p}=(1+x^p)^2\) over \(\F_p\).
Every polynomial in the line \(\F_ph\) is therefore divisible
by \(P\).

Filteredness of \(b\) would put
\(c(m_p^{(1)})^2\) in \(F_{2p}M_2\). Reducing this assertion
after multiplication by \(p^2\) puts \(cf^2\) in \(\F_ph\).
Evaluating at the \(a\) in (A-26f) yields
\(cf(a)^2=0\), contrary to \(p\nmid c\) and \(f(a)\ne0\).
This proves that every such bilinear map is zero. The
desired associated-graded product has bottom coefficient
\(T^{[1]}T^{[1]}=2T^{[2]}\), so it is excluded already by
the choice \(p=3\).
\end{proof}

\begin{proof}[Proof of the all-space assertion]
We use the cofiber sequences of finite filtration stages
\[
 F_{w-1}B(X)\longrightarrow F_wB(X)
                 \longrightarrow \operatorname{gr}_wB(X).
 \tag{A-26h}
\]
The weightwise PD comparison identifies, for a connected \(X\),
\[
 H^q(\operatorname{gr}_wB(X))=
 \begin{cases}H^q(X;\Z),&0\le q\le w,\\
 0,&q>w,\end{cases}
 \tag{A-26i}
\]
with the standard tail coefficient action and weighted cup
product. Negative cohomology is zero.
Every weight-\(w\), degree-zero class in (A-26i) lifts through
the coefficient map \(R_{\le w}\to F_wB(X)\), by compatibility
in (A-26a). Its connecting homomorphism in (A-26h) is
therefore zero.

Take a finite simplicial circle \(X=S^1\). Induction on \(w\)
in (A-26h) shows that \(F_wB(S^1)\) has no cohomology above
degree one or below degree zero, that its \(H^0\) is
\(R_{\le w}\), and that
\[
 0\longrightarrow H^1F_{w-1}B(S^1)
  \longrightarrow H^1F_wB(S^1)
  \longrightarrow
  \begin{cases}\Z,&w\ge1,\\0,&w=0\end{cases}
  \longrightarrow0.
 \tag{A-26j}
\]
In particular all the maps on \(H^1\) are injective.
Exactness of filtered colimits now gives a filtered \(R\)-module
\[
 M_1=H^1B(S^1),\qquad \operatorname{gr}M_1\cong I_1.
 \tag{A-26k}
\]
Here the filtration is the image of \(H^1F_wB(S^1)\);
the injectivity just proved identifies it with that group.
The graded module action is the required \(S\)-action by
the multiplicative compatibility in (A-26a).

Now let \(Y=S^1\times S^1\), with its two projections.
For each \(w\ge1\), the two classes of
\(H^1(\operatorname{gr}_wB(Y))\cong\Z^2\) lift from
\(H^1F_wB(S^1)\) by (A-26j) and naturality of the projections.
Their connecting maps to \(H^2F_{w-1}B(Y)\) in (A-26h)
are consequently zero. The degree-zero connecting maps
vanish as before, and there is no associated-graded cohomology
above degree two. Finite-stage induction gives
\[
 0\longrightarrow H^2F_{w-1}B(Y)
  \longrightarrow H^2F_wB(Y)
  \longrightarrow
  \begin{cases}\Z,&w\ge2,\\0,&w<2\end{cases}
  \longrightarrow0.
 \tag{A-26l}
\]
It follows, again using exact filtered colimits, that
\[
 M_2=H^2B(Y),\qquad \operatorname{gr}M_2\cong I_2
\]
as filtered \(R\)-modules with their induced graded action.
No convergence assertion about an infinite spectral sequence
has been used.

Pull back an element of \(M_1\) along the first projection
and another along the second, and take their cup product.
This defines a filtered \(R\)-bilinear pairing
\[
 M_1\otimes_R M_1\longrightarrow M_2.
\]
Choose the orientation of the torus so that the two ordinary
degree-one generators have product its degree-two generator.
The associated-graded pairing is then precisely
\(I_1\otimes_S I_1\to I_2\). Its value on the bottom generators
has coefficient \(2\), so it is nonzero. This contradicts the
first assertion of the theorem.
\end{proof}

\paragraph{An exhaustive finite certificate without a lattice classification.}
The independent program
\texttt{code/D\_filtered\_lattice\_gate.py} enumerates all
\(3\)-local filtered polynomial lattices of the required type
through weight five for \(M_1\) and weight six for \(M_2\).
Its completeness is elementary. In polynomial degree \(d\),
the lattice contains the coefficient basis vector
\(\beta_d=\binom Td\), and its prescribed leading-coefficient
ideal is \(i!/(d+i)!\) times \(\Z_{(3)}\).
Relative to the old degree-\((d-1)\) lattice and
\(\Z_{(3)}\beta_d\), the new index has \(3\)-part
\(3^e\), where
\[
 e=v_3\binom{d+i}{i}.
\]
If \(v_0,\ldots,v_{d-1}\) are the existing filtered basis,
every possible extension is generated by it and one vector
\[
 \frac{1}{3^e}
       \left(\beta_d+\sum_{j<d}a_jv_j\right),
       \qquad 0\le a_j<3^e.
 \tag{A-26m}
\]
Changing a coefficient by a multiple of \(3^e\) adds an
old lattice element; the unit part of the index only
rescales the new generator. Thus (A-26m) lists every
possible local lattice, including all lower coefficients.

The program tests every coefficient multiplication whose
degree stays within the cutoff, using the exact formula
\[
 \beta_a\beta_b=
 \sum_{c=\max(a,b)}^{a+b}
 \frac{c!}{(c-a)!(c-b)!(a+b-c)!}\,\beta_c.
\]
Triangular rational elimination tests membership by the
nonnegativity of the \(3\)-adic valuations of its coefficients.
The numbers of surviving lattices by polynomial degree are
\((1,1,9,3,3)\) for \(M_1\) and \((1,3,3,3,3)\) for \(M_2\).
Exactly three lattices of each type survive; none of their
nine pairs allows all required products of total weight
at most six. The corrected output
\texttt{logs/D\_filtered\_lattice\_gate.json} records their
bases and all counts, with 2005 exact membership checks.
This finite certificate independently checks the \(c=2\)
obstruction. The proof above establishes the stronger
all-prime, all-filtration statement.
